\chapter{Numerical Linear Algebra}

Every formula in this book has been exact. We wrote $\inv{\mat{A}}$, we wrote
$\sqrt{2}$, we wrote $1/3$, and we trusted those symbols to mean precisely what they
say. A real computer does not get that luxury. It stores numbers in a fixed number of
bits, so $1/3$ becomes $0.3333333\ldots$ chopped off somewhere, $\sqrt{2}$ becomes a
nearby fraction, and every arithmetic step nudges the answer by a hair. Most of the
time those hairs are harmless. Sometimes -- and this is the whole subject of the
chapter -- they conspire, amplify, and hand you an answer that is confidently,
catastrophically wrong.

Numerical linear algebra is the study of how to do linear algebra on a machine
\emph{and trust the result}. It splits into two questions that we will keep carefully
apart:

\begin{itemize}
  \item \textbf{Conditioning} is a property of the \emph{problem}. Some problems are
    inherently delicate: tiny changes to the input force large changes in the true
    answer. No algorithm, however clever, can rescue an ill-conditioned problem,
    because the input itself is already uncertain.
  \item \textbf{Stability} is a property of the \emph{algorithm}. A stable algorithm
    does not add much trouble of its own; an unstable one can wreck even a
    well-conditioned problem through its own internal rounding.
\end{itemize}

Hold those two words apart and most of the confusion in this field evaporates. The
condition number measures the problem. Stability measures the method. You want both:
a well-conditioned problem solved by a stable algorithm.

We will build the floating-point picture, define the condition number geometrically,
see why Gaussian elimination needs pivoting and why QR is the safer workhorse, then
turn to iterative methods -- Jacobi, Gauss--Seidel, and the conjugate gradient idea --
that solve enormous sparse systems without ever forming an inverse. We close with the
power method, the simplest eigenvalue algorithm, and a page of practical warnings you
will be glad to have.

\section{Floating point: numbers with a budget}

A computer stores a real number the way scientific notation does, but in binary and
with a fixed budget of digits. A double-precision (``float64'') number looks like
\[
  \pm\, 1.\underbrace{b_1 b_2 \cdots b_{52}}_{\text{52 bits}} \times 2^{e},
\]
a sign, a $52$-bit fraction (the \emph{mantissa}), and an exponent $e$. Because the
mantissa has only $52$ bits, the gap between one representable number and the next is
relative, not absolute: near $1$ the spacing is about $2^{-52}\approx 2.2\times
10^{-16}$, and near $1000$ it is a thousand times larger. The single most important
constant in the subject is this spacing at $1$, called \textbf{machine epsilon}:
\[
  \eps_{\text{mach}} \;=\; 2^{-52} \;\approx\; 2.22\times 10^{-16}.
\]

\begin{intuitionbox}
Think of the real line as a ruler whose tick marks get \emph{farther apart} as you
move away from zero, but always stay the same distance apart \emph{relative} to where
you are -- about $16$ significant decimal digits everywhere. When you store a number,
you snap it to the nearest tick. The snap costs at most half a tick, so the relative
error of storing any number $x$ is bounded:
\[
  \frac{\abs{\text{fl}(x) - x}}{\abs{x}} \;\le\; \tfrac12\,\eps_{\text{mach}}.
\]
Here $\text{fl}(x)$ means ``$x$ rounded to the nearest float.'' Sixteen digits sounds
like plenty. The danger is never one rounding -- it is what happens when millions of
them interact.
\end{intuitionbox}

\begin{figure}[H]
\centering
\begin{tikzpicture}[scale=1.0,>=Stealth]
  \draw[thick] (0,0)--(11,0);
  % evenly spaced ticks near 1 (dense), widening as we move right
  \foreach \x in {0.0,0.4,0.8,1.2,1.6,2.0,2.4} {
    \draw[vecblue] (\x,0.0)--(\x,0.18);
  }
  \foreach \x in {3.2,4.0,4.8,5.6,6.4} {
    \draw[vecblue] (\x,0.0)--(\x,0.18);
  }
  \foreach \x in {7.6,9.2,10.8} {
    \draw[vecblue] (\x,0.0)--(\x,0.18);
  }
  \node[below] at (1.2,-0.1){\small near $1$: ticks tight};
  \node[below] at (5.0,-0.1){\small near $10$};
  \node[below] at (9.2,-0.1){\small near $100$: ticks wide};
  \draw[<->,trapred,thick] (0.0,0.45)--(0.4,0.45);
  \node[trapred,above] at (0.2,0.5){\scriptsize $\eps_{\text{mach}}$};
  \draw[<->,trapred,thick] (7.6,0.45)--(9.2,0.45);
  \node[trapred,above] at (8.4,0.5){\scriptsize $\sim 100\,\eps_{\text{mach}}$};
\end{tikzpicture}
\caption{Floating-point numbers are a ruler whose tick marks spread apart as you move
from zero, but stay a constant \emph{relative} distance apart -- about $16$ significant
digits everywhere. Storing a real number snaps it to the nearest tick, costing at most
half a gap. The absolute gap near $100$ is about $100$ times the gap near $1$.}
\label{ch33:ruler}
\end{figure}

The IEEE~754 standard guarantees that the four basic operations are \emph{correctly
rounded}: the computed $x \oplus y$ equals the true $x+y$ rounded to the nearest
float. Symbolically, every operation introduces a tiny relative slip,
\[
  \text{fl}(x \circ y) \;=\; (x \circ y)(1+\delta), \qquad \abs{\delta}\le \tfrac12\eps_{\text{mach}},
\]
for each of $\circ \in \{+,-,\times,\div\}$. That little $(1+\delta)$ is the atom of
all rounding analysis. Stack billions of them and you want them to stay small and
\emph{not} reinforce.

\subsection{The cardinal sin: catastrophic cancellation}

Multiplication and division are gentle: they perturb the relative error by one
$\delta$ and move on. Addition is gentle too, \emph{unless you subtract two nearly
equal numbers}. That is where digits die.

\begin{example}[Subtracting close numbers]
\label{ch33:cancel}
Suppose a machine keeps only $6$ significant decimal digits and you compute
\[
  a = 1.23456\underline{7}, \qquad b = 1.23456\underline{1},
\]
each already rounded to $6$ digits (so the underlined seventh digit is gone). Their
true difference, had you kept full precision, might be anything from $0.0000050$ to
$0.0000070$ -- you simply do not know, because the digits that distinguish them were
thrown away. The stored result $a-b = 0.000006 = 6.00000\times 10^{-6}$ looks like it
has six good digits, but in truth it has \emph{one}, maybe zero. The leading digits
agreed and cancelled, promoting the formerly-tiny rounding errors into the leading
position.
\end{example}

This is \textbf{catastrophic cancellation}, and it is the villain behind most
numerical disasters. The subtraction itself is computed correctly -- the problem is
that it exposes errors that were already lurking in $a$ and $b$. The lesson, which we
will see echoed in elimination and least squares: \emph{avoid subtracting quantities
that are nearly equal whenever an algebraically equivalent route avoids it.}

\begin{labbox}
The quadratic formula $x = \dfrac{-b \pm \sqrt{b^2-4ac}}{2a}$ hides this trap. Take
$a=1$, $b=-200$, $c=1$. Then $\sqrt{b^2-4ac}=\sqrt{40000-4}=\sqrt{39996}\approx
199.98999\ldots$. The ``$+$'' root, $\tfrac{200+199.98999}{2}\approx 199.995$, is fine.
But the ``$-$'' root computes $\tfrac{200-199.98999}{2}$ -- a subtraction of two nearly
equal numbers -- and loses most of its digits. The cure is to compute the safe root
first, then get the other from $x_1 x_2 = c/a$, here $x_2 = 1/x_1$. Try both ways on a
calculator that shows few digits and watch the difference.
\end{labbox}

\begin{trapbox}
``Double precision has $16$ digits, so I have $16$ good digits.'' No. You \emph{start}
with $16$ digits. Each catastrophic cancellation can erase a clump of them at once. A
computation that subtracts two numbers agreeing in their first $12$ digits keeps about
$4$. The bits are not lost in the subtraction -- they were never there to begin with;
the subtraction just removes the disguise.
\end{trapbox}

\section{Conditioning: how delicate is the problem?}

Forget computers for a moment. Suppose someone hands you a linear system
$\mat{A}\vec{x}=\vec{b}$ and warns you that $\vec{b}$ is only known to within a little
wiggle $\Delta\vec{b}$ -- a measurement error, say. How much can the solution
$\vec{x}$ wiggle in response? If a small $\Delta\vec{b}$ can force a large
$\Delta\vec{x}$, the problem is \textbf{ill-conditioned}, and \emph{no} algorithm can
save you: the answer genuinely is uncertain.

Geometrically, solving $\mat{A}\vec{x}=\vec{b}$ in $\R^2$ means intersecting two
lines (the two equations). If the lines cross briskly at a wide angle, nudging one
line barely moves the crossing point. But if the two lines are \emph{nearly parallel},
the intersection point is a long, skinny wedge: a hair's nudge to either line slides
the crossing point a huge distance (Figure~\ref{ch33:nearparallel}). That near-parallel
picture is exactly what ill-conditioning looks like.

\begin{figure}[H]
\centering
\begin{tikzpicture}[scale=0.95,>=Stealth]
  % Left: well-conditioned, wide angle
  \begin{scope}[xshift=0cm]
    \draw[->,gray!50] (-0.3,0)--(3.2,0);
    \draw[->,gray!50] (0,-0.3)--(0,3.2);
    \draw[very thick,vecblue] (-0.2,2.6)--(3.0,0.2);
    \draw[very thick,vecred]  (0.1,-0.2)--(2.6,3.0);
    \fill[black] (1.45,1.4) circle (1.6pt);
    \node at (1.5,3.4){\small wide angle: stable};
    \node at (1.5,-0.7){\small (well-conditioned)};
  \end{scope}
  % Right: ill-conditioned, near parallel
  \begin{scope}[xshift=5.3cm]
    \draw[->,gray!50] (-0.3,0)--(3.2,0);
    \draw[->,gray!50] (0,-0.3)--(0,3.2);
    \draw[very thick,vecblue] (-0.2,1.0)--(3.0,2.5);
    \draw[very thick,vecred]  (-0.2,1.25)--(3.0,2.65);
    \fill[black] (1.4,1.75) circle (1.6pt);
    \draw[->,trapred,thick] (1.4,1.75)--(2.55,2.3);
    \node[trapred] at (2.9,2.0){\small slides far};
    \node at (1.5,3.4){\small near-parallel: fragile};
    \node at (1.5,-0.7){\small (ill-conditioned)};
  \end{scope}
\end{tikzpicture}
\caption{Two equations are two lines. When they cross at a wide angle (left), the
solution is robust to perturbing either line. When they are nearly parallel (right),
a tiny change to a line slides the intersection a long way -- the hallmark of an
ill-conditioned system.}
\label{ch33:nearparallel}
\end{figure}

\subsection{The condition number, made precise}

Let us put a number on that delicacy. Start from $\mat{A}\vec{x}=\vec{b}$ and perturb
the right side to $\vec{b}+\Delta\vec{b}$, giving a new solution
$\vec{x}+\Delta\vec{x}$. Subtracting the two systems, $\mat{A}\,\Delta\vec{x} =
\Delta\vec{b}$, so $\Delta\vec{x} = \inv{\mat{A}}\,\Delta\vec{b}$. Taking norms and
using the operator-norm inequality $\norm{\mat{M}\vec{v}}\le\norm{\mat{M}}\norm{\vec{v}}$,
\[
  \norm{\Delta\vec{x}} \le \norm{\inv{\mat{A}}}\,\norm{\Delta\vec{b}},
  \qquad\text{and from } \mat{A}\vec{x}=\vec{b},\quad
  \norm{\vec{b}} \le \norm{\mat{A}}\,\norm{\vec{x}}.
\]
Multiply the two inequalities together and rearrange to get the relative error bound:
\[
  \boxed{\;\frac{\norm{\Delta\vec{x}}}{\norm{\vec{x}}}
   \;\le\; \underbrace{\norm{\mat{A}}\,\norm{\inv{\mat{A}}}}_{\displaystyle \kappa(\mat{A})}
   \;\cdot\; \frac{\norm{\Delta\vec{b}}}{\norm{\vec{b}}}\;.}
\]
The amplification factor in the middle is the \textbf{condition number}
$\kappa(\mat{A}) = \norm{\mat{A}}\,\norm{\inv{\mat{A}}}$. It tells you how a relative
error in the input can blow up into a relative error in the output.

\begin{formalbox}
For the $2$-norm (the usual Euclidean length), the matrix norm $\norm{\mat{A}}$ equals
the largest singular value $\sigma_{\max}$, and $\norm{\inv{\mat{A}}} = 1/\sigma_{\min}$.
So
\[
  \kappa_2(\mat{A}) \;=\; \frac{\sigma_{\max}}{\sigma_{\min}},
\]
the ratio of the largest to the smallest singular value of $\mat{A}$. A condition
number near $1$ is ideal (orthogonal matrices have $\kappa_2 = 1$ exactly). A condition
number of $10^{k}$ means you can expect to lose about $k$ decimal digits of accuracy:
if $\kappa_2(\mat{A})\approx 10^{8}$ and you start with $16$ digits, trust only about
$8$.
\end{formalbox}

\begin{connectionbox}
The condition number is the SVD wearing a geometric hat. Recall from the SVD chapter
that $\mat{A}$ maps the unit sphere into an ellipsoid whose semi-axes are the singular
values $\sigma_1\ge\cdots\ge\sigma_n$. The condition number $\kappa_2(\mat{A}) = \sigma_{\max}/\sigma_{\min}$
is simply the \emph{aspect ratio} of that ellipsoid (Figure~\ref{ch33:condellipse}). A round
sphere (aspect ratio near $1$) means $\mat{A}$ stretches every direction about equally, and
reversing it is benign. But when $\sigma_{\max} \gg \sigma_{\min}$, $\mat{A}$ squashes the unit sphere into an elongated needle. When you solve $\mat{A}\vec{x}=\vec{b}$, the inverse $\inv{\mat{A}}$ must stretch that needle back into a round sphere -- which means any microscopic disturbance or floating-point wobble along the needle's skinny waist is magnified by the enormous factor $1/\sigma_{\min}$. Conditioning, geometry, and the SVD are one unified idea.
\end{connectionbox}

\begin{figure}[H]
\centering
\begin{tikzpicture}[scale=0.88,>=Stealth]
  % Panel 1: Round input / well-conditioned
  \begin{scope}[shift={(0,0)}]
    \node[anchor=south,font=\footnotesize\scshape,text=spacegray!90!black] at (0,2.2) {1. Well-Conditioned ($\kappa \approx 1.5$)};
    \draw[fill=codebg,draw=spacegray!30,rounded corners=4pt] (-2.2,-2.1) rectangle (2.2,2.0);
    \draw[gray!40,dashed] (0,0) circle (1.0);
    \draw[fill=sourceblue!15,draw=sourceblue!80,line width=1.3pt] (0,0) ellipse (1.3 and 0.85);
    \draw[->,vecblue,thick] (0,0) -- (1.3,0) node[below=1pt,font=\scriptsize]{$\sigma_{\max}$};
    \draw[->,vecred,thick] (0,0) -- (0,0.85) node[left=1pt,font=\scriptsize]{$\sigma_{\min}$};
    \node[font=\scriptsize,align=center,text=spacegray!90!black] at (0,-1.35) {Round image: $\mat{A}$ preserves\\relative proportions};
  \end{scope}

  % Panel 2: The Needle (severe squashing)
  \begin{scope}[shift={(5.2,0)}]
    \node[anchor=south,font=\footnotesize\scshape,text=spacegray!90!black] at (0,2.2) {2. Ill-Conditioned Needle ($\kappa \gg 1$)};
    \draw[fill=codebg,draw=spacegray!30,rounded corners=4pt] (-2.2,-2.1) rectangle (2.2,2.0);
    \draw[gray!40,dashed] (0,0) circle (1.0);
    \draw[fill=quotientorange!15,draw=quotientorange!90,line width=1.3pt] (0,0) ellipse (1.8 and 0.16);
    \draw[->,vecblue,thick] (0,0) -- (1.8,0) node[below right=-2pt,font=\scriptsize]{$\sigma_{\max}$};
    \draw[->,vecred,thick] (0,0) -- (0,0.16) node[above=2pt,font=\scriptsize]{$\sigma_{\min} \approx 0$};
    \node[font=\scriptsize,align=center,text=spacegray!90!black] at (0,-1.35) {Flattened needle: information\\along $\sigma_{\min}$ is almost lost};
  \end{scope}

  % Panel 3: Inversion and Error Magnification
  \begin{scope}[shift={(10.4,0)}]
    \node[anchor=south,font=\footnotesize\scshape,text=spacegray!90!black] at (0,2.2) {3. Inversion Error Blow-Up};
    \draw[fill=codebg,draw=spacegray!30,rounded corners=4pt] (-2.2,-2.1) rectangle (2.2,2.0);
    
    % tiny input disturbance circle Delta b
    \draw[fill=gray!20,draw=gray!60] (-1.0,0.7) circle (0.2);
    \node[font=\tiny,above] at (-1.0,0.95) {$\Delta\vec{b}$ (tiny noise)};

    % Arrow showing inversion stretch
    \draw[->,thick,vecred] (-0.6,0.7) to[bend left=15] node[midway,above,font=\scriptsize]{$\inv{\mat{A}}$} (0.3,0.7);

    % huge stretched error ellipse Delta x
    \draw[fill=vecred!15,draw=vecred!80,line width=1.2pt] (1.1,-0.1) ellipse (0.25 and 1.35);
    \draw[<->,vecred!90!black,thick] (1.5,-1.35) -- (1.5,1.25)
      node[midway,right,font=\scriptsize]{$\norm{\Delta\vec{x}} \propto \frac{\norm{\Delta\vec{b}}}{\sigma_{\min}}$};
    \node[font=\scriptsize,align=center,text=spacegray!90!black] at (0,-1.5) {Tiny input error $\Delta\vec{b}$ expands\\into massive solution error $\Delta\vec{x}$};
  \end{scope}
\end{tikzpicture}
\caption{The condition number as geometric distortion and error magnification. Left: When $\kappa_2(\mat{A}) = \sigma_{\max}/\sigma_{\min}$ is small, $\mat{A}$ deforms the unit sphere into a nearly round ellipse, making inversion stable. Middle: When $\kappa_2(\mat{A})$ is large, $\mat{A}$ squashes the unit sphere into an elongated needle. Right: Inverting the system $\mat{A}\vec{x}=\vec{b}$ requires $\inv{\mat{A}}$ to expand the needle back out, magnifying any tiny perturbation $\Delta\vec{b}$ perpendicular to the needle's axis by $1/\sigma_{\min} = \kappa_2(\mat{A})/\sigma_{\max}$.}
\label{ch33:condellipse}
\end{figure}

\begin{example}[A mildly ill-conditioned $2\times 2$]
\label{ch33:cond2x2}
Take
\[
  \mat{A} = \begin{bmatrix} 1 & 1 \\ 1 & 1.001 \end{bmatrix}.
\]
The two rows are \emph{almost} the same -- two nearly parallel lines. Its determinant
is $1\cdot 1.001 - 1\cdot 1 = 0.001$, tiny, so the inverse is large:
\[
  \inv{\mat{A}} = \frac{1}{0.001}\begin{bmatrix} 1.001 & -1 \\ -1 & 1 \end{bmatrix}
   = \begin{bmatrix} 1001 & -1000 \\ -1000 & 1000 \end{bmatrix}.
\]
Now watch the amplification directly. With $\vec{b}=\colvec{2\\2.001}$ the exact
solution is $\vec{x}=\colvec{1\\1}$. Perturb the second entry of $\vec{b}$ by just
$0.001$ to $\vec{b}'=\colvec{2\\2.002}$:
\[
  \vec{x}' = \inv{\mat{A}}\vec{b}' = \colvec{1001\cdot 2 - 1000\cdot 2.002 \\ -1000\cdot 2 + 1000\cdot 2.002}
   = \colvec{0\\2}.
\]
A $0.05\%$ change in $\vec{b}$ flipped the answer from $(1,1)$ to $(0,2)$ -- a $100\%$
change in $\vec{x}$. The condition number here is roughly $\kappa_2(\mat{A})\approx
4000$, and indeed the relative error was amplified by a few thousand -- right up against
that worst-case bound.
\end{example}

\begin{trapbox}
A small determinant does \emph{not} mean ill-conditioned, and a large one does not
mean well-conditioned -- the determinant has the wrong units. Scale $\mat{A}$ by
$\tfrac{1}{10}$ and its determinant drops by $10^{n}$, yet $\kappa(\mat{A})$ does not
change at all, because both $\norm{\mat{A}}$ and $\norm{\inv{\mat{A}}}$ rescale and the
factors cancel. Conditioning is about the \emph{ratio} of stretches
$\sigma_{\max}/\sigma_{\min}$, not the overall volume $\det$. Judge sensitivity by
$\kappa$, never by $\det$.
\end{trapbox}

\begin{walkbox}
The Hilbert matrix has entries $H_{ij}=1/(i+j-1)$ -- innocent-looking fractions like
$1, \tfrac12, \tfrac13,\ldots$ on the diagonals. Its condition number explodes: about
$500$ for $3\times 3$, roughly half a million for $5\times 5$, about $1.6\times 10^{13}$
by $10\times 10$, and past $10^{16}$ -- beyond all double-precision hope -- a few sizes
later. On a walk, ask yourself why ``smooth, slowly varying'' matrices like this -- rows
that look nearly proportional -- should be the dangerous ones, while a matrix of random
$\pm 1$ entries is usually tame. (Hint: think about how parallel the rows are.)
\end{walkbox}

\section{Stability: elimination versus QR}

The condition number is about the problem. Now we ask about the \emph{algorithm}. Even
a perfectly well-conditioned system can be wrecked by a careless method. The classic
example is Gaussian elimination without pivoting.

\begin{historybox}
In 1947, as the electronic computing era dawned with the room-sized ENIAC running thousands of glowing vacuum tubes, John von Neumann and Herman Goldstine published a landmark paper: \emph{Numerical Inverting of Matrices of High Order}. Before electronic computers, solving linear systems was done by human teams with mechanical desk calculators; a $10 \times 10$ matrix took weeks of grueling hand labor. The ENIAC promised to solve $100 \times 100$ systems in minutes.

Yet a deep sense of dread gripped numerical analysts. Leading mathematicians, including Norbert Wiener and Alan Turing, feared that Gaussian elimination on digital computers would be rendered useless by cascading roundoff errors. They worried that tiny rounding slips would double at each step, growing exponentially ($2^n$) and obliterating every significant digit after just twenty or thirty rows. Von Neumann and Goldstine introduced the matrix condition number $\kappa(\mat{A}) = \norm{\mat{A}} \norm{\mat{A}\inv}$ to bound error growth, but early estimates remained pessimistic.

The breakthrough came in the 1960s from James H. Wilkinson, Turing's former collaborator at Britain's National Physical Laboratory. Wilkinson invented \textbf{backward error analysis}. Instead of asking how far the computer's answer $\hat{\vec{x}}$ was from the true solution $\vec{x}$ (which is hopelessly magnified if the matrix is sensitive), Wilkinson turned the question around: \emph{what slightly perturbed problem did the machine solve exactly?} He proved that Gaussian elimination with partial pivoting solves $(\mat{A} + \Delta\mat{A})\hat{\vec{x}} = \vec{b}$ where the perturbation $\Delta\mat{A}$ is on the order of machine precision. Elimination does not create errors; it merely reveals the sensitivity that was already baked into the problem.

At the same time, the Cold War inspired another computational revolution. In 1965, James Cooley and John Tukey published the Fast Fourier Transform (FFT). Tukey was serving on President Kennedy's Science Advisory Committee, searching for ways to detect underground Soviet nuclear tests in seismic sensor data. Slashing the cost of Fourier transforms from $\mathcal{O}(n^2)$ to $\mathcal{O}(n \log n)$ made real-time signal processing possible on digital hardware---an algorithm Carl Friedrich Gauss had already derived in 1805 to track asteroid orbits, only to leave it unpublished in a desk drawer.
\end{historybox}

\subsection{Why elimination needs pivoting}

Recall elimination subtracts a multiple of one row from another to create zeros. The
multiplier is $m = (\text{entry to kill})/(\text{pivot})$. If the pivot is tiny, $m$ is
enormous, and multiplying a row by an enormous number, then subtracting, drowns the
original information in roundoff.

\begin{example}[A tiny pivot ruins a benign system]
\label{ch33:pivot}
Solve, on a machine carrying only $3$ significant digits,
\[
  \begin{bmatrix} 0.001 & 1 \\ 1 & 1 \end{bmatrix}
  \begin{bmatrix} x_1 \\ x_2 \end{bmatrix} =
  \begin{bmatrix} 1 \\ 2 \end{bmatrix}.
\]
This system is perfectly well-conditioned; the true answer is about
$x_1 = 1.001\ldots,\ x_2 = 0.999\ldots$, i.e. close to $(1,1)$.

\textbf{Without pivoting}, the pivot is $0.001$ and the multiplier is
$m = 1/0.001 = 1000$. Eliminating the second row:
\[
  \text{row}_2 \leftarrow \text{row}_2 - 1000\cdot\text{row}_1:
  \quad (1 - 1000\cdot 0.001)\,x_1 + (1 - 1000\cdot 1)\,x_2 = 2 - 1000\cdot 1.
\]
The new $x_2$ coefficient is $1 - 1000 = -999$, and the right side is $2-1000=-998$.
But with only $3$ digits, the original ``$1$'' in $1-1000$ is utterly swamped: the
machine stores $-1000$ and $-1000$, losing the very entries that carried the answer.
It then reports $x_2 = (-1000)/(-1000) = 1.00$, and back-substitution gives
$x_1 = (1 - 1\cdot 1)/0.001 = 0$. The computed answer is $(0, 1)$ -- the first
coordinate is completely wrong, off by $100\%$, on a problem that was never hard.

\textbf{With partial pivoting}, you first \emph{swap} so the largest-magnitude entry in
the column, the $1$, sits in the pivot:
\[
  \begin{bmatrix} 1 & 1 \\ 0.001 & 1 \end{bmatrix}
  \begin{bmatrix} x_1 \\ x_2 \end{bmatrix} =
  \begin{bmatrix} 2 \\ 1 \end{bmatrix}.
\]
Now the multiplier is $0.001$, small and harmless. Eliminating:
$x_2$ coefficient becomes $1 - 0.001 = 0.999$, right side $1 - 0.002 = 0.998$, so
$x_2 = 0.998/0.999 \approx 0.999$, and back-substitution gives $x_1 \approx 1.00$. The
answer $(1.00, 0.999)$ is correct to the precision available. Same problem, same
arithmetic precision -- only the order of operations changed.
\end{example}

\begin{focusbox}
\textbf{Partial pivoting:} before eliminating in a column, swap rows so the
largest-magnitude available entry becomes the pivot. This keeps every multiplier
$\abs{m}\le 1$, so no row is blown up before being subtracted. It is cheap (a row swap),
and it is why \texttt{A\textbackslash b} in MATLAB and \texttt{numpy.linalg.solve} are
reliable in practice. Plain LU \emph{without} pivoting is numerically reckless; LU
\emph{with} partial pivoting, $\mat{P}\mat{A}=\mat{L}\mat{U}$, is the standard.
\end{focusbox}

\subsection{Backward stability, in one sentence}

Pivoted Gaussian elimination is \textbf{backward stable}: the answer it returns is the
\emph{exact} solution of a slightly perturbed system $(\mat{A}+\mat{E})\hat{\vec{x}} =
\vec{b}$, where $\mat{E}$ is tiny relative to $\mat{A}$. That is the gold standard for a
numerical algorithm. You did not solve your problem exactly; you solved a problem
exactly that sits within rounding distance of yours. Whether the \emph{answer} is then
accurate depends on the conditioning -- a stable algorithm on an ill-conditioned
problem still gives a large error, but it is the problem's fault, not the method's. The
two ideas combine into the rule of thumb every numerical analyst carries:
\[
  \text{forward error} \;\lesssim\; \kappa(\mat{A}) \times \text{backward error}.
\]

\subsection{Why QR is the safe workhorse}

There is a class of matrices with condition number exactly $1$: orthogonal matrices
$\mat{Q}$, with $\mat{Q}\T\mat{Q}=\mat{I}$. They are rotations and reflections; they
preserve length, $\norm{\mat{Q}\vec{x}}=\norm{\vec{x}}$, so they cannot amplify any
error. This is the deep reason orthogonal transformations are the safest tools in the
numerical toolbox: \emph{they never make a vector's relative error worse.}

The QR factorization $\mat{A}=\mat{Q}\mat{R}$ (Gram--Schmidt's stable cousins, the
Householder and Givens methods, build it) exploits this. To solve
$\mat{A}\vec{x}=\vec{b}$ you write $\mat{Q}\mat{R}\vec{x}=\vec{b}$, multiply by
$\mat{Q}\T$ -- free, because it is orthogonal and harmless -- to get
$\mat{R}\vec{x}=\mat{Q}\T\vec{b}$, then back-substitute through the triangular $\mat{R}$.
The orthogonal step transports the data without inflating any error
(Figure~\ref{ch33:orthogonal}).

\begin{figure}[H]
\centering
\begin{tikzpicture}[scale=1.0,>=Stealth]
  \draw[gray!50] (-0.4,0)--(3.2,0);
  \draw[gray!50] (0,-0.4)--(0,3.0);
  % original disk (error blob) -> rotated, same size
  \fill[vecblue!25] (0.9,0.7) ellipse (0.45 and 0.45);
  \node[vecblue] at (0.9,1.45){\small error blob};
  \draw[->,thick] (1.7,0.7)--(2.55,0.7);
  \node at (2.1,1.0){\small $\mat{Q}\T$};
  \fill[targetgreen!25] (3.3,0.7) ellipse (0.45 and 0.45);
  \node at (5.0,0.7){\small same radius: $\kappa(\mat{Q})=1$};
\end{tikzpicture}
\caption{An orthogonal matrix rotates the error blob but never stretches it. Because
$\norm{\mat{Q}\vec{x}}=\norm{\vec{x}}$, applying $\mat{Q}\T$ leaves relative errors
untouched -- the reason QR-based solvers are so robust.}
\label{ch33:orthogonal}
\end{figure}

\begin{strangbox}
Strang teaches Gaussian elimination first because it is how you compute by hand and how
$\mat{A}=\mat{L}\mat{U}$ is born. He then introduces $\mat{A}=\mat{Q}\mat{R}$ for
least squares. The numerical punchline, which a first pass through Strang can leave
implicit: for solving square systems, pivoted LU is fast and stable enough; but for
\emph{least squares} $\min\norm{\mat{A}\vec{x}-\vec{b}}$, never form the normal
equations $\mat{A}\T\mat{A}\vec{x}=\mat{A}\T\vec{b}$ on a computer, because
$\kappa(\mat{A}\T\mat{A}) = \kappa(\mat{A})^2$ -- you square the condition number and
can lose twice as many digits. Solve via QR (or the SVD) instead. That single warning
is the practical heart of this chapter.
\end{strangbox}

\begin{applicationbox}
\textbf{Why your machine-learning library uses QR/SVD under the hood.} A linear
regression with $p$ features fits $\min_{\vec\beta}\norm{\mat{X}\vec\beta - \vec{y}}$.
The textbook ``closed form'' $\vec\beta = (\mat{X}\T\mat{X})\inv\mat{X}\T\vec{y}$ is
mathematically correct and numerically treacherous: when features are correlated (which
real features always are), $\mat{X}$ is ill-conditioned and $\mat{X}\T\mat{X}$ doubles
the damage. That is why \texttt{scikit-learn}'s \texttt{LinearRegression} calls
LAPACK's SVD-based least-squares routine, not the normal-equation formula, and why
\emph{ridge regression} -- adding $\lambda\mat{I}$ to $\mat{X}\T\mat{X}$ -- works so
well: it lifts the smallest singular values away from zero, slashing the condition
number and stabilizing the fit. Regularization is, in part, a conditioning fix.
\end{applicationbox}

\section{Iterative methods: solving without inverting}

Everything so far has been \emph{direct}: a finite recipe (elimination, QR) that
produces the answer after a fixed number of steps, roughly $\tfrac23 n^3$ for an
$n\times n$ system. For $n=10^6$ -- a million unknowns, routine in scientific computing
and graphics -- that is $10^{18}$ operations and a dense $10^{12}$-entry matrix that
would not fit in memory. Direct methods die there.

But the giant matrices of practice are almost always \textbf{sparse}: mostly zeros, a
handful of nonzeros per row, because each unknown couples only to its few neighbors (a
pixel to adjacent pixels, a node to adjacent nodes). For these we use
\textbf{iterative methods}, which never modify $\mat{A}$ and never form $\inv{\mat{A}}$.
They only ever multiply $\mat{A}$ by a vector -- cheap when $\mat{A}$ is sparse -- and
they sneak up on the answer, generating a sequence $\vec{x}^{(0)},
\vec{x}^{(1)},\vec{x}^{(2)},\ldots$ that converges to $\vec{x}$.

\subsection{Jacobi: everyone updates from last round}

Split $\mat{A}$ into its diagonal $\mat{D}$ and the rest $\mat{R}=\mat{A}-\mat{D}$.
The system $\mat{A}\vec{x}=\vec{b}$ becomes $\mat{D}\vec{x} = \vec{b}-\mat{R}\vec{x}$,
and if we already had $\vec{x}$ on the right we could read off $\vec{x}$ on the left.
We do not, so we iterate: plug in the current guess on the right to get the next guess
on the left,
\[
  \vec{x}^{(k+1)} = \inv{\mat{D}}\bigl(\vec{b} - \mat{R}\,\vec{x}^{(k)}\bigr).
\]
Component by component, the recipe is plain: solve the $i$-th equation for
$x_i$, treating every \emph{other} unknown as known from the previous round,
\[
  x_i^{(k+1)} = \frac{1}{a_{ii}}\Bigl(b_i - \sum_{j\ne i} a_{ij}\,x_j^{(k)}\Bigr).
\]

\begin{example}[Two Jacobi sweeps by hand]
\label{ch33:jacobi}
Solve
\[
  \begin{bmatrix} 4 & 1 \\ 2 & 5 \end{bmatrix}\vec{x} = \colvec{1\\7},
  \qquad\text{true answer } \vec{x}\approx\colvec{-0.1111 \\ 1.4444}.
\]
The Jacobi updates are $x_1 \leftarrow (1 - x_2)/4$ and $x_2 \leftarrow (7-2x_1)/5$.
Start at $\vec{x}^{(0)}=(0,0)$:
\[
  \begin{aligned}
  x_1^{(1)} &= (1-0)/4 = 0.25, & x_2^{(1)} &= (7-0)/5 = 1.4,\\
  x_1^{(2)} &= (1-1.4)/4 = -0.1, & x_2^{(2)} &= (7-2\cdot 0.25)/5 = 1.3,\\
  x_1^{(3)} &= (1-1.3)/4 = -0.075, & x_2^{(3)} &= (7-2\cdot(-0.1))/5 = 1.44.
  \end{aligned}
\]
Already after three sweeps we are at $(-0.075,\,1.44)$, marching toward
$(-0.11,\,1.44)$. Each sweep used only the previous sweep's numbers -- so all updates
could run in parallel.
\end{example}

\subsection{Gauss-Seidel: use the news as soon as you hear it}

Jacobi wastes information: while computing the new $x_2$ it still uses the \emph{stale}
$x_1$, even though a fresh $x_1$ was just computed this very sweep. Gauss--Seidel fixes
that -- use each new component the instant it is available:
\[
  x_i^{(k+1)} = \frac{1}{a_{ii}}\Bigl(b_i - \sum_{j<i} a_{ij}\,x_j^{(k+1)}
     - \sum_{j>i} a_{ij}\,x_j^{(k)}\Bigr).
\]
The only change from Jacobi is the superscript on the first sum: $j<i$ uses the
already-updated values $\vec{x}^{(k+1)}$. Redo the example: $x_1^{(1)}=(1-0)/4=0.25$ as
before, but now $x_2^{(1)}=(7-2\cdot 0.25)/5 = 1.3$ immediately uses the fresh $0.25$.
Next sweep $x_1^{(2)}=(1-1.3)/4=-0.075$, then $x_2^{(2)}=(7-2\cdot(-0.075))/5=1.43$.
After two sweeps Gauss--Seidel sits at $(-0.075,1.43)$, where Jacobi needed three. It
typically converges about twice as fast, at the cost of being inherently sequential
(each component waits for the one before it).

\begin{intuitionbox}
Both methods are ``relax toward agreement.'' Picture each equation as a constraint, and
each sweep as nudging every unknown to satisfy its own equation given everyone else's
current value. Like a roomful of people each adjusting to their neighbors, the whole
configuration settles toward the consensus that satisfies all equations at once.
Gauss--Seidel just lets each person hear the latest gossip before adjusting, so the
room settles faster.
\end{intuitionbox}

\begin{formalbox}
These iterations converge for \emph{any} starting guess precisely when the
\textbf{spectral radius} of the iteration matrix is below $1$. For Jacobi the iteration
matrix is $\mat{M}=-\inv{\mat{D}}\mat{R}$, and $\vec{x}^{(k+1)}-\vec{x} =
\mat{M}(\vec{x}^{(k)}-\vec{x})$, so the error is multiplied by $\mat{M}$ each step and
dies iff every eigenvalue of $\mat{M}$ has $\abs{\lambda}<1$. A clean sufficient
condition you can check by eye: \textbf{strict diagonal dominance},
$\abs{a_{ii}} > \sum_{j\ne i}\abs{a_{ij}}$ for every row -- each diagonal entry bigger
than the rest of its row combined -- guarantees both Jacobi and Gauss--Seidel converge.
Our $\left[\begin{smallmatrix}4&1\\2&5\end{smallmatrix}\right]$ qualifies: $4>1$ and
$5>2$.
\end{formalbox}

\begin{trapbox}
Iterative methods are not guaranteed to converge. If $\mat{A}$ is not diagonally
dominant (or, more generally, if the iteration matrix has spectral radius $\ge 1$),
Jacobi and Gauss--Seidel can stall or \emph{diverge to infinity}. Always check
diagonal dominance, or a related structural property like positive-definiteness, before
trusting them -- and watch the residual $\norm{\mat{A}\vec{x}^{(k)}-\vec{b}}$ actually
shrinking.
\end{trapbox}

\subsection{Conjugate gradient: the heavyweight}

For the most important case -- $\mat{A}$ symmetric positive definite -- there is a
method that leaves Jacobi and Gauss--Seidel far behind: the \textbf{conjugate gradient}
(CG) method. The geometry behind it is what makes it work.

When $\mat{A}$ is symmetric positive definite, solving $\mat{A}\vec{x}=\vec{b}$ is
\emph{exactly} the same as minimizing the bowl-shaped quadratic
\[
  \phi(\vec{x}) = \tfrac12\,\vec{x}\T\mat{A}\vec{x} - \vec{b}\T\vec{x}.
\]
Its gradient is $\nabla\phi = \mat{A}\vec{x}-\vec{b}$, which is zero exactly at the
solution. So solving the linear system \emph{is} rolling to the bottom of a
paraboloid bowl. Because $\mat{A}$ is positive definite, the bowl really is a bowl
(curving up in every direction), and its level sets are nested ellipsoids -- stretched
into a long valley when $\mat{A}$ is ill-conditioned.

\begin{figure}[H]
\centering
\begin{tikzpicture}[scale=1.0,>=Stealth]
  % elongated ellipse contours (a stretched valley)
  \foreach \r in {0.5,1.0,1.5,2.0} {
    \draw[gray!55] (0,0) ellipse ({1.7*\r} and {0.6*\r});
  }
  \fill[black] (0,0) circle (1.6pt);
  \node[below] at (0,-0.15){\small minimum $=\vec{x}$};
  % steepest descent: zig-zag
  \coordinate (s0) at (-3.0,1.0);
  \coordinate (s1) at (-1.4,-0.5);
  \coordinate (s2) at (-0.9,0.45);
  \coordinate (s3) at (-0.35,-0.18);
  \draw[->,trapred,thick] (s0)--(s1);
  \draw[->,trapred,thick] (s1)--(s2);
  \draw[->,trapred,thick] (s2)--(s3);
  \node[trapred] at (-2.6,1.35){\small steepest descent};
  % CG: two clean steps
  \coordinate (c0) at (2.9,1.0);
  \coordinate (c1) at (1.4,-1.0);
  \draw[->,vecblue,very thick] (c0)--(c1);
  \draw[->,vecblue,very thick] (c1)--(0,0);
  \node[vecblue] at (2.6,1.35){\small conjugate gradient};
\end{tikzpicture}
\caption{Solving $\mat{A}\vec{x}=\vec{b}$ for symmetric positive definite $\mat{A}$
means finding the bottom of an ellipsoidal bowl. Naive steepest descent (red)
zig-zags across the stretched valley, taking tiny corrective steps forever. Conjugate
gradient (blue) chooses $\mat{A}$-orthogonal directions and never undoes its own
progress, reaching the bottom of an $n$-dimensional bowl in at most $n$ steps.}
\label{ch33:cg}
\end{figure}

The obvious way down is \emph{steepest descent}: always step in the direction of
steepest decrease (the negative gradient). On a round bowl that goes straight to the
bottom, but on a long ellipsoidal valley it zig-zags maddeningly, each step ruining the
progress of the last (Figure~\ref{ch33:cg}, red). Conjugate gradient's insight is to
pick directions that are ``$\mat{A}$-orthogonal'' (conjugate): $\vec{p}_i\T\mat{A}\,
\vec{p}_j=0$ for $i\ne j$. With such directions, a step that fixes one direction never
spoils another, so there is no zig-zag and no backtracking. In exact arithmetic CG
reaches the exact solution of an $n\times n$ system in at most $n$ steps -- and usually
far fewer, because it converges fast once it has explored the directions that matter.

\begin{intuitionbox}
Steepest descent is like walking downhill in a narrow canyon by always facing the
steepest local slope: you bounce wall to wall, creeping forward. Conjugate gradient is
like realizing the canyon has a few principal directions and dispatching each one
\emph{once and for all}, so you stride to the bottom in a handful of decisive moves.
The number of steps you really need is governed by $\sqrt{\kappa(\mat{A})}$ -- which is
why \emph{preconditioning} (multiplying by an approximate inverse to shrink $\kappa$)
is what makes CG practical on real problems.
\end{intuitionbox}

\begin{applicationbox}
Conjugate gradient and its relatives are the quiet engine of large-scale computation.
When a graphics engine simulates cloth or smoke, when an aircraft is stress-tested by
finite-element analysis, or when a modern recommendation system solves a sparse linear
system in millions of unknowns, a Krylov-subspace iterative solver (CG for symmetric
problems, GMRES for general ones) is almost always inside the loop. They win because
they touch $\mat{A}$ only through matrix-vector products, which respect sparsity and
parallelize across thousands of cores -- something a dense $\inv{\mat{A}}$ could never
do.
\end{applicationbox}

\section{The power method: finding the top eigenvalue}

Iterative ideas also conquer eigenvalues. The simplest eigenvalue algorithm, the
\textbf{power method}, computes the largest-magnitude eigenvalue and its eigenvector
using nothing but repeated matrix-vector multiplication.

The idea: start with almost any vector $\vec{v}^{(0)}$ and hit it with $\mat{A}$ over
and over. Write $\vec{v}^{(0)}$ in an eigenbasis,
$\vec{v}^{(0)} = c_1\vec{x}_1 + c_2\vec{x}_2 + \cdots + c_n\vec{x}_n$, where
$\abs{\lambda_1} > \abs{\lambda_2}\ge\cdots$. Then
\[
  \mat{A}^k\vec{v}^{(0)} = c_1\lambda_1^k\vec{x}_1 + c_2\lambda_2^k\vec{x}_2 + \cdots
   = \lambda_1^k\Bigl(c_1\vec{x}_1 + c_2(\tfrac{\lambda_2}{\lambda_1})^k\vec{x}_2 + \cdots\Bigr).
\]
Every ratio $\lambda_j/\lambda_1$ has magnitude below $1$, so its $k$-th power decays to
zero. The dominant eigendirection $\vec{x}_1$ survives; everything else fades. After
enough multiplications the vector points (almost) exactly along $\vec{x}_1$. To keep the
numbers from overflowing or underflowing, normalize after each step:
\[
  \vec{w} = \mat{A}\,\vec{v}^{(k)}, \qquad
  \vec{v}^{(k+1)} = \frac{\vec{w}}{\norm{\vec{w}}}.
\]
The eigenvalue itself falls out as the Rayleigh quotient
$\lambda_1 \approx \vec{v}^{(k)\mathsf T}\mat{A}\,\vec{v}^{(k)}$, or simply as the
factor by which the vector grew before normalizing.

\begin{example}[Power method in three steps]
\label{ch33:power}
Take
\[
  \mat{A} = \begin{bmatrix} 2 & 1 \\ 1 & 2 \end{bmatrix},
\]
whose true eigenvalues are $3$ (eigenvector $(1,1)$) and $1$ (eigenvector $(1,-1)$).
Start with the deliberately lopsided $\vec{v}^{(0)}=\colvec{1\\0}$ and just track the
raw vector $\mat{A}^k\vec{v}^{(0)}$:
\[
  \mat{A}\colvec{1\\0}=\colvec{2\\1},\quad
  \mat{A}\colvec{2\\1}=\colvec{5\\4},\quad
  \mat{A}\colvec{5\\4}=\colvec{14\\13},\quad
  \mat{A}\colvec{14\\13}=\colvec{41\\40}.
\]
Watch the ratio of the two components head for $1$ (the direction of $(1,1)$):
$2/1$, then $5/4$, then $14/13$, then $41/40$. And the growth factor per step heads
for $3$: from $\colvec{14\\13}$ to $\colvec{41\\40}$ the first component went
$14\to 41$, a factor $2.93$, closing on the true dominant eigenvalue $\lambda_1=3$.
Four multiplications and we already have two good digits
(Figure~\ref{ch33:powerfig}).
\end{example}

\begin{figure}[H]
\centering
\begin{tikzpicture}[scale=1.1,>=Stealth]
  \draw[gray!50] (-0.3,0)--(3.0,0);
  \draw[gray!50] (0,-0.3)--(0,3.0);
  \draw[dashed,structurepurple] (0,0)--(2.6,2.6) node[right]{\small $\vec{x}_1$ dir $(1,1)$};
  \draw[->,vecblue,thick] (0,0)--(2.0,0.0) node[below]{\small $v^{(0)}$};
  \draw[->,vecblue,thick] (0,0)--(1.79,0.89);
  \draw[->,vecblue,thick] (0,0)--(1.56,1.25);
  \draw[->,vecblue,thick] (0,0)--(1.41,1.34);
  \node[vecblue] at (1.9,1.7){\small swings toward $\vec{x}_1$};
\end{tikzpicture}
\caption{The power method in action on $\left[\begin{smallmatrix}2&1\\1&2\end{smallmatrix}\right]$.
Each multiplication by $\mat{A}$ swings the iterate closer to the dominant eigendirection
$(1,1)$; the subdominant component shrinks by the factor $\lambda_2/\lambda_1 = 1/3$
each step.}
\label{ch33:powerfig}
\end{figure}

\begin{applicationbox}
\textbf{PageRank is the power method.} Google's original ranking treated the web as a
giant matrix $\mat{A}$ where $A_{ij}$ is the probability a random surfer hops from page
$j$ to page $i$. The importance of every page is the dominant eigenvector of that
billion-by-billion matrix -- and you find it by exactly the iteration above: start with
a uniform vector, multiply by $\mat{A}$, normalize, repeat. Because $\mat{A}$ is
enormously sparse (each page links to only a handful of others), each multiplication is
cheap, and the dominant eigenvalue is conveniently $1$ with a clear gap to the next, so
the power method converges in a few dozen sweeps over the entire web. The whole of
linear algebra's spectral theory, run iteratively, ranked the internet.
\end{applicationbox}

\begin{remark}
The power method only finds the \emph{top} eigenvalue, and it stalls if the top two are
equal in magnitude (no gap, nothing to win the race). To get other eigenvalues you
shift and invert (apply the power method to $(\mat{A}-\mu\mat{I})\inv$ to find the
eigenvalue nearest $\mu$), or deflate, or -- for all eigenvalues at once -- use the
\textbf{QR algorithm}, the true workhorse: repeatedly factor $\mat{A}_k=\mat{Q}_k\mat{R}_k$
and reassemble $\mat{A}_{k+1}=\mat{R}_k\mat{Q}_k$. Astonishingly, this similarity
iteration drives the matrix toward upper-triangular form, reading the eigenvalues off
the diagonal. It is the power method's grown-up sibling and what \texttt{eig} actually
calls.
\end{remark}

\section{Practical warnings to carry with you}

Let us collect the field's hard-won wisdom into a checklist you can actually use.

\begin{focusbox}
\textbf{Numerical commandments.}
\begin{enumerate}
  \item \textbf{Never invert a matrix to solve a system.} To solve $\mat{A}\vec{x}=
    \vec{b}$, factor and back-substitute (\texttt{solve}, \texttt{A\textbackslash b}) --
    do not compute $\inv{\mat{A}}$ and multiply. Forming $\inv{\mat{A}}$ is slower, less
    accurate, and destroys sparsity.
  \item \textbf{Never solve least squares via the normal equations.} Use QR or SVD;
    $\mat{A}\T\mat{A}$ squares the condition number.
  \item \textbf{Check the condition number} (\texttt{cond(A)}) before trusting a
    result. If $\kappa\approx 10^{k}$, distrust the last $k$ digits.
  \item \textbf{Avoid subtracting nearly equal numbers.} Reformulate to dodge
    catastrophic cancellation when you can.
  \item \textbf{Never test floating-point numbers for exact equality.} Write
    $\abs{a-b} < \text{tol}$, not $a == b$; rounding makes exact equality a lie.
  \item \textbf{Exploit structure.} Sparse, symmetric, positive definite, banded,
    triangular -- each unlocks a faster, more stable specialized routine. Tell your
    library what you know.
  \item \textbf{Watch the residual, and remember conditioning.} A small residual
    $\norm{\mat{A}\hat{\vec{x}}-\vec{b}}$ does not by itself guarantee a small error
    $\norm{\hat{\vec{x}}-\vec{x}}$ -- the gap between them is the condition number.
\end{enumerate}
\end{focusbox}

\begin{aestheticbox}
Pure linear algebra hands us perfect objects --
$\inv{\mat{A}}$, exact eigenvalues, clean subspaces. Numerical linear algebra asks the
humbler, harder question: \emph{which of these can survive contact with a finite
machine?} The answers have a pleasing shape. Orthogonal matrices survive because they preserve
length. Sparse problems survive because we never form the dense inverse. Ill-conditioned
problems warn us through $\kappa$, and stable algorithms promise to add no
lies of their own. The condition number and the SVD turn out to be the same idea seen
twice. The whole field is a meditation on doing infinite-precision mathematics with
finite means -- and on knowing, precisely, how much to trust the answer.
\end{aestheticbox}

\begin{recoverybox}
Before the exercises, settle the two pillars. \emph{Conditioning} is about the problem:
$\kappa(\mat{A})=\sigma_{\max}/\sigma_{\min}$ measures how much a relative error in the
data can be amplified into the answer; nothing can fix an ill-conditioned problem.
\emph{Stability} is about the algorithm: a backward-stable method (pivoted LU, QR)
solves a nearby problem exactly and adds little of its own. The master inequality
$\text{forward error}\lesssim\kappa\times\text{backward error}$ binds them. If those two
sentences feel solid, you have the chapter.
\end{recoverybox}

\subsection*{What We Can Now Do (Capabilities Unlocked)}
\begin{itemize}
  \item \textbf{Rigorous Error Bounding:} We can calculate the condition number $\kappa(\mat{A}) = \sigma_{\max}/\sigma_{\min}$ and establish strict theoretical bounds on solution error via the Master Inequality: $\frac{\norm{\hat{\vec{x}}-\vec{x}}}{\norm{\vec{x}}} \le \frac{\kappa(\mat{A})}{1 - \kappa(\mat{A})\frac{\norm{\Delta\mat{A}}}{\norm{\mat{A}}}} \left(\frac{\norm{\Delta\mat{A}}}{\norm{\mat{A}}} + \frac{\norm{\Delta\vec{b}}}{\norm{\vec{b}}}\right)$.
  \item \textbf{Safe Factorizations with Partial Pivoting:} By enforcing row swaps so that the pivot is always the largest entry in the column ($|a_{kk}| \ge |a_{ik}|$), we guarantee multipliers $|m_{ik}| \le 1$, preventing pivot explosion and ensuring backward stability.
  \item \textbf{Orthogonal Stability via QR:} Using Householder reflections or Givens rotations, we perform QR factorization where $\mat{Q}$ is perfectly isometric ($\kappa_2(\mat{Q}) = 1$), eliminating condition number squaring in least-squares problems ($\mat{A}\T\mat{A}$).
  \item \textbf{Sparse Iterative Scaling:} With Jacobi, Gauss--Seidel, and the Power Method, we can solve systems and extract dominant eigenvectors for sparse matrices with millions of rows in $\mathcal{O}(n)$ per iteration, without ever materializing an explicit dense inverse.
\end{itemize}

\subsection*{Honest Downsides and Limitations}
\begin{itemize}
  \item \textbf{Cubic Scaling Bottleneck:} Direct solvers (LU, QR, Cholesky, SVD) scale as $\mathcal{O}(n^3)$ in time and $\mathcal{O}(n^2)$ in memory. For modern problems where $n = 10^7$ (such as Web graph ranking or neural networks), direct factorizations would require exabytes of RAM and centuries of compute.
  \item \textbf{The Curse of Ill-Conditioning:} Backward stability guarantees that the computed answer is exact for a nearby problem; it cannot rescue a problem where $\kappa(\mat{A}) > 1/\eps_{\text{mach}} \approx 10^{16}$. In such degenerate cases, double precision loses all 16 digits, returning numerical noise.
  \item \textbf{Slow Iterative Convergence under Weak Spectral Gaps:} The Power Method and basic relaxation schemes converge at a geometric rate $r = |\lambda_2/\lambda_1|$. When $\lambda_2 \approx \lambda_1$, the error ratio approaches 1, requiring tens of thousands of matrix-vector products.
  \item \textbf{Memory Thrashing and Fill-in:} Performing Gaussian elimination on a sparse matrix often destroys sparsity, creating dense ``fill-in'' that exhausts computer memory unless sophisticated graph reordering (such as Approximate Minimum Degree) is applied.
\end{itemize}

\subsection*{Where We Go Next in the Book}
\begin{itemize}
  \item \textbf{Chapter 34 (Linear Algebra of Machine Learning):} When models have billions of parameters, storing full matrices or computing $\mathcal{O}(n^3)$ factorizations is impossible. We discover how machine learning embraces low-rank matrix decompositions (LoRA), stochastic gradient descent, and dimension reduction to train massive transformers.
  \item \textbf{Chapter 35 (Graphs, Networks, and Markov Chains):} We apply the Power Method to the multi-billion-node Google PageRank problem, discovering how the Perron--Frobenius theorem guarantees convergence.
  \item \textbf{Chapter 36 (Optimization and Quadratic Forms):} We replace slow stationary iterations (Jacobi/Gauss-Seidel) with the celebrated Conjugate Gradient method, which accelerates convergence on symmetric positive-definite systems to $\mathcal{O}(\sqrt{\kappa})$.
\end{itemize}

\section{Exercises}

\begin{exercise}[Machine epsilon]
Double precision has a $52$-bit mantissa, so $\eps_{\text{mach}}=2^{-52}$. (a) Estimate
$\eps_{\text{mach}}$ in decimal and confirm it gives roughly $16$ significant digits.
(b) Explain in one sentence why the smallest number you can add to $1.0$ and get a
result different from $1.0$ is about $\eps_{\text{mach}}$, while near $1000$ that
smallest difference is about $1000\,\eps_{\text{mach}}$.
\end{exercise}

\begin{exercise}[Catastrophic cancellation]
Working to $4$ significant decimal digits throughout, compute $f(x)=\dfrac{1-\cos x}
{x^2}$ at $x=0.001$ directly. The true value is very close to $1/2$. (a) What do you
get, and why is it badly wrong? (b) Using $1-\cos x = 2\sin^2(x/2)$, recompute and
compare. (c) State which step in the direct route caused the loss.
\end{exercise}

\begin{exercise}[Condition number of a $2\times 2$]
Let $\mat{A}=\left[\begin{smallmatrix}1 & 0\\ 0 & \delta\end{smallmatrix}\right]$ for a
small $\delta>0$. (a) Find both singular values and write down $\kappa_2(\mat{A})$. (b)
For $\delta=10^{-6}$, how many decimal digits of accuracy do you expect to lose when
solving $\mat{A}\vec{x}=\vec{b}$? (c) Sketch the image of the unit circle under
$\mat{A}$ and relate its shape to $\kappa$.
\end{exercise}

\begin{exercise}[Determinant is not conditioning]
Let $\mat{A}=10^{-3}\mat{I}_3$ (the $3\times 3$ identity times $0.001$). (a) Compute
$\det\mat{A}$ -- it is tiny. (b) Compute $\kappa_2(\mat{A})$. (c) Explain in one or two
sentences why the small determinant is no cause for alarm here, tying it back to the
trap that conditioning is about a ratio of stretches, not a volume.
\end{exercise}

\begin{exercise}[Pivoting saves the day]
Consider $\left[\begin{smallmatrix}10^{-4} & 1\\ 1 & 1\end{smallmatrix}\right]\vec{x}=
\colvec{1\\2}$, solved on a machine with $3$ significant digits. (a) Solve by Gaussian
elimination \emph{without} pivoting and report the (wrong) answer. (b) Swap rows
(partial pivoting) and solve again. (c) State the multiplier in each case and explain
which one caused the trouble.
\end{exercise}

\begin{exercise}[Diagonal dominance and convergence]
For $\mat{A}=\left[\begin{smallmatrix}5 & 2\\ 1 & 4\end{smallmatrix}\right]$ and
$\vec{b}=\colvec{7\\6}$: (a) Verify strict diagonal dominance. (b) Carry out three
Jacobi sweeps from $\vec{x}^{(0)}=(0,0)$. (c) Carry out three Gauss--Seidel sweeps from
the same start, and confirm it is closer to the true solution $(\tfrac{16}{18},
\tfrac{23}{18})\approx(0.889,1.278)$ after the same number of sweeps.
\end{exercise}

\begin{exercise}[Why not the normal equations]
Suppose $\mat{A}$ (tall, full column rank) has $\kappa_2(\mat{A})=10^{6}$. (a) What is
$\kappa_2(\mat{A}\T\mat{A})$? (b) Roughly how many decimal digits would you expect to
lose solving the least-squares problem through the normal equations versus through QR,
starting from double precision? (c) In one sentence, state the practical rule this
implies.
\end{exercise}

\begin{exercise}[Power method by hand]
Let $\mat{A}=\left[\begin{smallmatrix}3 & 1\\ 0 & 2\end{smallmatrix}\right]$, whose
eigenvalues are $3$ and $2$. Starting from $\vec{v}^{(0)}=\colvec{1\\1}$, perform three
power-method steps (you may skip normalization and just track the vector and the ratio
of components). (a) Toward which eigenvector does the iterate swing? (b) Estimate the
dominant eigenvalue from the growth factor. (c) The subdominant component dies by what
factor each step, and what does that predict about the speed of convergence?
\end{exercise}

\begin{exercise}[Conjugate gradient geometry]
Take the symmetric positive definite $\mat{A}=\left[\begin{smallmatrix}1 & 0\\ 0 &
100\end{smallmatrix}\right]$ and $\vec{b}=\colvec{1\\1}$. (a) Write the quadratic
$\phi(\vec{x})$ and describe the shape of its level sets. (b) Explain why steepest
descent from a generic start zig-zags badly, referencing the condition number
$\kappa_2(\mat{A})=100$. (c) Argue that conjugate gradient finishes in at most $2$
steps here, and say what general statement this illustrates.
\end{exercise}

\begin{exercise}[A proof: orthogonal matrices are perfectly conditioned]
Let $\mat{Q}$ be a real $n\times n$ orthogonal matrix, $\mat{Q}\T\mat{Q}=\mat{I}$.
(a) Show $\norm{\mat{Q}\vec{x}}_2=\norm{\vec{x}}_2$ for every $\vec{x}$, hence
$\norm{\mat{Q}}_2=1$. (b) Show $\inv{\mat{Q}}=\mat{Q}\T$ is also orthogonal, so
$\norm{\inv{\mat{Q}}}_2=1$. (c) Conclude $\kappa_2(\mat{Q})=1$, and explain in words why
this makes orthogonal transformations the safest building blocks in numerical linear
algebra.
\end{exercise}

\begin{exercise}[Residual versus error -- a challenge]
Using $\mat{A}=\left[\begin{smallmatrix}1 & 1\\ 1 & 1.0001\end{smallmatrix}\right]$ and
the true solution $\vec{x}=\colvec{1\\1}$ (so $\vec{b}=\colvec{2\\2.0001}$), consider
the approximate solution $\hat{\vec{x}}=\colvec{2\\0}$. (a) Compute the residual
$\vec{r}=\vec{b}-\mat{A}\hat{\vec{x}}$ and its norm. (b) Compute the error
$\hat{\vec{x}}-\vec{x}$ and its norm. (c) The residual is small but the error is large.
Compute $\kappa_2(\mat{A})$ approximately and explain how the master inequality
$\text{forward error}\lesssim\kappa\times\text{backward error}$ accounts for the gap.
\end{exercise}
